% =========================================================================
% BÀI 35: CÁC NGÀY TRONG TUẦN
% TOÀN BỘ CÂU HỎI ĐƯỢC VẼ HÌNH MINH HỌA TRỰC QUAN 100% BẰNG TIKZ 2D
% =========================================================================

% --- ĐỊNH NGHĨA MÀU SẮC AN TOÀN ---
\colorlet{myteal}{cyan!70!blue}
\colorlet{myorange}{orange}
\colorlet{mypurple}{purple}
\colorlet{myblue}{blue}
\colorlet{myred}{red}

% --- CÁC HÀM VẼ HÌNH TRỰC QUAN ĐẶC TRƯNG CHO BÀI 35 ---

% 1. Vẽ tờ lịch bóc tường / để bàn
\providecommand{\VeToLich}[4]{% #1=(x,y), #2=Thứ, #3=Số/Nội dung, #4=màu
  \begin{scope}[shift={(#1)}]
    \fill[black!70] (-0.45,0.7) circle (0.06cm);
    \fill[black!70] (0.45,0.7) circle (0.06cm);
    \draw[line width=1.1pt, black!50] (-0.45,0.7) -- (-0.45,0.86);
    \draw[line width=1.1pt, black!50] (0.45,0.7) -- (0.45,0.86);
    \draw[fill=white, draw=#4!85, line width=1.1pt, rounded corners=3pt] (-1.1,-0.7) rectangle (1.1,0.7);
    \draw[fill=#4!85, draw=#4!85, rounded corners=2pt] (-1.1,0.24) rectangle (1.1,0.7);
    \node[font=\bfseries\tiny, color=white] at (0,0.47) {#2};
    \node[font=\bfseries\scriptsize, color=#4!90!black] at (0,-0.23) {#3};
  \end{scope}
}

% 2. Dải băng 7 ngày trong tuần
\providecommand{\VeDaiBayNgay}[2]{% #1=(x,y), #2=ngày nổi bật (1..7, 0 là không)
  \begin{scope}[shift={(#1)}, scale=0.58]
    \foreach \i/\thu/\mau in {
      1/{Thứ 2}/myblue,
      2/{Thứ 3}/myblue,
      3/{Thứ 4}/myblue,
      4/{Thứ 5}/myblue,
      5/{Thứ 6}/myblue,
      6/{Thứ 7}/myorange,
      7/{Chủ nhật}/myred
    } {
      \pgfmathsetmacro{\x}{(\i-4)*1.92}
      \ifnum\i=#2
        \draw[fill=\mau!25, draw=\mau!90, line width=1.3pt, rounded corners=3pt] (\x-0.82,-0.55) rectangle (\x+0.82,0.55);
        \node[font=\bfseries\footnotesize, color=\mau!90!black] at (\x,0) {\thu};
        \node[above, font=\bfseries\tiny, color=red!85] at (\x,0.6) {$\blacktriangledown$};
      \else
        \draw[fill=white, draw=gray!70, line width=0.7pt, rounded corners=3pt] (\x-0.82,-0.55) rectangle (\x+0.82,0.55);
        \node[font=\bfseries\scriptsize, color=black!85] at (\x,0) {\thu};
      \fi
    }
  \end{scope}
}

% 3. Mũi tên thời gian giữa 2 sự kiện
\providecommand{\VeMuiTenLich}[3]{% #1=(x,y), #2=độ dài, #3=nhãn
  \begin{scope}[shift={(#1)}]
    \draw[->, >=stealth, line width=1.1pt, orange!90!black] (0,0) -- (#2,0);
    \node[above=1pt, font=\bfseries\tiny, color=blue!90!black] at (#2*0.5, 0) {#3};
  \end{scope}
}

% =========================================================================
% BẮT ĐẦU NỘI DUNG 20 CÂU HỎI
% =========================================================================

\textbf{Phần 1. Trắc nghiệm lựa chọn}
\setcounterpageref{ex}{0}

% --- TRANG 1: CÂU 1 ĐẾN CÂU 4 ---

% CÂU 1: Một tuần lễ có 7 ngày
\begin{ex} %[1H9N35]
Dải các ngày trong một tuần lễ được xếp theo thứ tự:
\begin{center}\vspace{-2.5mm}
\begin{tikzpicture}
  \VeDaiBayNgay{0,0}{0}
\end{tikzpicture}
\vspace{-2.5mm}\end{center}
Một tuần lễ có tất cả bao nhiêu ngày?
\begin{multicols}{4}
\choice
{$5$ ngày}
{$6$ ngày}
{\True $7$ ngày}
{$8$ ngày}
\end{multicols}
\loigiai{
Một tuần lễ có đúng $7$ ngày: Thứ Hai, Thứ Ba, Thứ Tư, Thứ Năm, Thứ Sáu, Thứ Bảy và Chủ nhật.\\
\textbf{Đáp án: C.}
}
\end{ex}

% CÂU 2: Ngày đầu tiên trong tuần đi học là Thứ Hai
\begin{ex} %[1H9N35]
Tờ lịch đầu tuần đón các em học sinh đến trường:
\begin{center}\vspace{-2.5mm}
\begin{tikzpicture}
  \VeToLich{0,0}{THỨ HAI}{Tuần mới}{myblue}
\end{tikzpicture}
\vspace{-2.5mm}\end{center}
Ngày đầu tiên trong tuần lễ đi học là ngày nào?
\begin{multicols}{4}
\choice
{Chủ nhật}
{\True Thứ Hai}
{Thứ Ba}
{Thứ Bảy}
\end{multicols}
\loigiai{
Ngày đầu tiên trong tuần lễ đi học theo quy ước là ngày Thứ Hai.\\
\textbf{Đáp án: B.}
}
\end{ex}

% CÂU 3: Hai ngày nghỉ cuối tuần
\begin{ex} %[1H9H35]
Hai tờ lịch ngày nghỉ cuối tuần của các bạn nhỏ:
\begin{center}\vspace{-2.5mm}
\begin{tikzpicture}
  \VeToLich{-1.5,0}{THỨ BẢY}{Nghỉ học}{myorange}
  \VeToLich{1.5,0}{CHỦ NHẬT}{Nghỉ học}{myred}
\end{tikzpicture}
\vspace{-2.5mm}\end{center}
Học sinh tiểu học thường được nghỉ học vào những ngày nào trong tuần?
\begin{multicols}{2}
\choice
{\True Thứ Bảy và Chủ nhật}
{Thứ Hai và Thứ Ba}
{Thứ Năm và Thứ Sáu}
{Thứ Tư và Thứ Năm}
\end{multicols}
\loigiai{
Học sinh tiểu học thường được nghỉ học vào hai ngày cuối tuần là Thứ Bảy và Chủ nhật.\\
\textbf{Đáp án: A.}
}
\end{ex}

% CÂU 4: Số ngày đi học trong một tuần
\begin{ex} %[1H9H35]
Các ngày em đến trường học tập trong tuần:
\begin{center}\vspace{-2.5mm}
\begin{tikzpicture}[scale=0.68]
  \foreach \i/\thu in {1/{Thứ 2}, 2/{Thứ 3}, 3/{Thứ 4}, 4/{Thứ 5}, 5/{Thứ 6}} {
    \pgfmathsetmacro{\x}{(\i-3)*1.95}
    \draw[fill=blue!10, draw=blue!75, line width=0.8pt, rounded corners=3pt] (\x-0.85,-0.5) rectangle (\x+0.85,0.5);
    \node[font=\bfseries\scriptsize, color=blue!90!black] at (\x,0.12) {\thu};
    \node[font=\bfseries\tiny, color=black!65] at (\x,-0.25) {Đi học};
  }
\end{tikzpicture}
\vspace{-2.5mm}\end{center}
Trong một tuần lễ, em đi học bao nhiêu ngày (từ Thứ Hai đến Thứ Sáu)?
\begin{multicols}{4}
\choice
{\True $5$ ngày}
{$4$ ngày}
{$6$ ngày}
{$7$ ngày}
\end{multicols}
\loigiai{
Từ Thứ Hai đến Thứ Sáu có tất cả $5$ ngày đi học.\\
\textbf{Đáp án: A.}
}
\end{ex}

\newpage
% --- TRANG 2: CÂU 5 ĐẾN CÂU 8 ---

% CÂU 5: Hôm nay Thứ Ba -> Ngày mai là Thứ Tư
\begin{ex} %[1H9H35]
Quan sát mối quan hệ giữa Hôm nay và Ngày mai:
\begin{center}\vspace{-2.5mm}
\begin{tikzpicture}
  \VeToLich{-2.2,0}{THỨ BA}{Hôm nay}{myblue}
  \VeMuiTenLich{-0.8,0}{1.6}{Ngày mai}
  \draw[fill=yellow!25, draw=orange!80, line width=0.9pt, rounded corners=3pt] (1.1,-0.7) rectangle (3.3,0.7);
  \node[font=\bfseries\small, color=red!85] at (2.2,0.1) {$?$};
  \node[font=\bfseries\tiny, color=black!70] at (2.2,-0.35) {Ngày mai};
\end{tikzpicture}
\vspace{-2.5mm}\end{center}
Nếu hôm nay là Thứ Ba thì ngày mai là thứ mấy?
\begin{multicols}{4}
\choice
{Thứ Hai}
{Thứ Năm}
{Thứ Sáu}
{\True Thứ Tư}
\end{multicols}
\loigiai{
Ngày liền sau Thứ Ba là Thứ Tư. Vậy nếu hôm nay là Thứ Ba thì ngày mai là Thứ Tư.\\
\textbf{Đáp án: D.}
}
\end{ex}

% CÂU 6: Hôm nay Thứ Năm -> Hôm qua là Thứ Tư
\begin{ex} %[1H9H35]
Quan sát tờ lịch ngày Hôm nay:
\begin{center}\vspace{-2.5mm}
\begin{tikzpicture}
  \draw[fill=yellow!25, draw=orange!80, line width=0.9pt, rounded corners=3pt] (-3.3,-0.7) rectangle (-1.1,0.7);
  \node[font=\bfseries\small, color=red!85] at (-2.2,0.1) {$?$};
  \node[font=\bfseries\tiny, color=black!70] at (-2.2,-0.35) {Hôm qua};
  
  \VeMuiTenLich{-0.8,0}{1.6}{Tiếp theo}
  \VeToLich{2.2,0}{THỨ NĂM}{Hôm nay}{myblue}
\end{tikzpicture}
\vspace{-2.5mm}\end{center}
Nếu hôm nay là Thứ Năm thì hôm qua là thứ mấy?
\begin{multicols}{4}
\choice
{Thứ Ba}
{Thứ Sáu}
{\True Thứ Tư}
{Thứ Bảy}
\end{multicols}
\loigiai{
Ngày liền trước Thứ Năm là Thứ Tư. Vậy hôm qua là Thứ Tư.\\
\textbf{Đáp án: C.}
}
\end{ex}

% CÂU 7: Ngày liền trước Chủ nhật là Thứ Bảy
\begin{ex} %[1H9V35]
Hai ngày cuối tuần liên tiếp:
\begin{center}\vspace{-2.5mm}
\begin{tikzpicture}
  \draw[fill=yellow!25, draw=orange!80, line width=0.9pt, rounded corners=3pt] (-3.1,-0.7) rectangle (-0.9,0.7);
  \node[font=\bfseries\small, color=red!85] at (-2.0,0.1) {$?$};
  \node[font=\bfseries\tiny, color=black!70] at (-2.0,-0.35) {Liền trước};
  
  \VeMuiTenLich{-0.7,0}{1.4}{}
  \VeToLich{2.0,0}{CHỦ NHẬT}{Cuối tuần}{myred}
\end{tikzpicture}
\vspace{-2.5mm}\end{center}
Ngày đứng liền trước ngày Chủ nhật trong tuần lễ là ngày nào?
\begin{multicols}{4}
\choice
{Thứ Sáu}
{Thứ Năm}
{Thứ Hai}
{\True Thứ Bảy}
\end{multicols}
\loigiai{
Ngày đứng liền trước ngày Chủ nhật là ngày Thứ Bảy.\\
\textbf{Đáp án: D.}
}
\end{ex}

% CÂU 8: Hôm nay Thứ Sáu -> Ngày kia là Chủ nhật
\begin{ex} %[1H9V35]
Mối liên hệ giữa Hôm nay, Ngày mai và Ngày kia:
\begin{center}\vspace{-2.5mm}
\begin{tikzpicture}
  \VeToLich{-3.4,0}{THỨ SÁU}{Hôm nay}{myblue}
  \VeMuiTenLich{-2.0,0}{1.0}{Mai}
  \VeToLich{0,0}{THỨ BẢY}{Ngày mai}{myorange}
  \VeMuiTenLich{1.3,0}{1.0}{Kia}
  \draw[fill=yellow!25, draw=orange!80, line width=0.9pt, rounded corners=3pt] (2.6,-0.7) rectangle (4.8,0.7);
  \node[font=\bfseries\small, color=red!85] at (3.7,0.1) {$?$};
  \node[font=\bfseries\tiny, color=black!70] at (3.7,-0.35) {Ngày kia};
\end{tikzpicture}
\vspace{-2.5mm}\end{center}
Nếu hôm nay là Thứ Sáu thì ngày kia là ngày nào?
\begin{multicols}{4}
\choice
{Thứ Bảy}
{Thứ Hai}
{Thứ Tư}
{\True Chủ nhật}
\end{multicols}
\loigiai{
Hôm nay là Thứ Sáu $\rightarrow$ Ngày mai là Thứ Bảy $\rightarrow$ Ngày kia là Chủ nhật.\\
\textbf{Đáp án: D.}
}
\end{ex}

\newpage
% --- TRANG 3: CÂU 9 ĐẾN CÂU 12 ---

% CÂU 9: Kế hoạch tuần của Mai
\begin{ex} %[1H9V35]
Bảng kế hoạch tuần của bạn Mai:
\begin{center}\vspace{-2.5mm}
\begin{tikzpicture}
  \VeToLich{-3.4,0}{THỨ HAI}{Học bơi}{myblue}
  \VeToLich{0,0}{THỨ TƯ}{Học đàn}{myteal}
  \VeToLich{3.4,0}{THỨ BẢY}{Về quê}{myorange}
\end{tikzpicture}
\vspace{-2.5mm}\end{center}
Theo kế hoạch trên, bạn Mai về quê thăm ông bà vào ngày nào?
\begin{multicols}{4}
\choice
{\True Thứ Bảy}
{Thứ Hai}
{Thứ Tư}
{Chủ nhật}
\end{multicols}
\loigiai{
Quan sát tờ lịch ghi rõ: Thứ Bảy -- Về quê.\\
\textbf{Đáp án: A.}
}
\end{ex}

% CÂU 10: Ngày ở giữa Thứ Tư và Thứ Sáu là Thứ Năm
\begin{ex} %[1H9V35]
Đoàn tàu chở ba ngày liên tiếp trong tuần:
\begin{center}\vspace{-2.5mm}
\begin{tikzpicture}
  \draw[fill=blue!15, draw=blue!80, line width=0.8pt, rounded corners=3pt] (-3.4,-0.55) rectangle (-1.2,0.55);
  \node[font=\bfseries\footnotesize, color=blue!90!black] at (-2.3,0) {Thứ Tư};
  
  \draw[fill=yellow!25, draw=orange!80, line width=0.9pt, rounded corners=3pt] (-0.8,-0.55) rectangle (1.2,0.55);
  \node[font=\bfseries\large, color=red!85] at (0.2,0) {$?$};
  
  \draw[fill=blue!15, draw=blue!80, line width=0.8pt, rounded corners=3pt] (1.6,-0.55) rectangle (3.8,0.55);
  \node[font=\bfseries\footnotesize, color=blue!90!black] at (2.7,0) {Thứ Sáu};
\end{tikzpicture}
\vspace{-2.5mm}\end{center}
Ngày nằm ở giữa ngày Thứ Tư và ngày Thứ Sáu là ngày nào?
\begin{multicols}{4}
\choice
{Thứ Ba}
{\True Thứ Năm}
{Thứ Bảy}
{Chủ nhật}
\end{multicols}
\loigiai{
Thứ tự các ngày: Thứ Tư $\rightarrow$ Thứ Năm $\rightarrow$ Thứ Sáu. Ngày ở giữa là Thứ Năm.\\
\textbf{Đáp án: B.}
}
\end{ex}

% CÂU 11: Hôm nay Chủ nhật -> Hôm kia là Thứ Sáu
\begin{ex} %[1H9V35]
Tìm ngày trước đó $2$ ngày (Hôm kia):
\begin{center}\vspace{-2.5mm}
\begin{tikzpicture}
  \draw[fill=yellow!25, draw=orange!80, line width=0.9pt, rounded corners=3pt] (-4.6,-0.7) rectangle (-2.4,0.7);
  \node[font=\bfseries\small, color=red!85] at (-3.5,0.1) {$?$};
  \node[font=\bfseries\tiny, color=black!70] at (-3.5,-0.35) {Hôm kia};
  
  \draw[<-, >=stealth, line width=1.1pt, orange!90!black] (-2.2,0) -- (-1.3,0);
  \VeToLich{0,0}{THỨ BẢY}{Hôm qua}{myorange}
  \draw[<-, >=stealth, line width=1.1pt, orange!90!black] (1.3,0) -- (2.2,0);
  \VeToLich{3.5,0}{CHỦ NHẬT}{Hôm nay}{myred}
\end{tikzpicture}
\vspace{-2.5mm}\end{center}
Nếu hôm nay là Chủ nhật thì hôm kia là ngày nào?
\begin{multicols}{4}
\choice
{Thứ Bảy}
{\True Thứ Sáu}
{Thứ Năm}
{Thứ Tư}
\end{multicols}
\loigiai{
Hôm nay là Chủ nhật $\rightarrow$ Hôm qua là Thứ Bảy $\rightarrow$ Hôm kia là Thứ Sáu.\\
\textbf{Đáp án: B.}
}
\end{ex}

% CÂU 12: Sau Chủ nhật vòng lại Thứ Hai
\begin{ex} %[1H9C35]
Vòng tuần hoàn khép kín của các ngày trong tuần:
\begin{center}\vspace{-2.5mm}
\begin{tikzpicture}
  \VeToLich{-2.6,0}{CHỦ NHẬT}{Hết tuần}{myred}
  \draw[->, >=stealth, line width=1.2pt, red!85] (-1.2,0) -- (1.2,0);
  \node[above=1pt, font=\bfseries\tiny, color=purple!85] at (0,0) {Tuần mới bắt đầu};
  \draw[fill=yellow!25, draw=orange!80, line width=0.9pt, rounded corners=3pt] (1.5,-0.7) rectangle (3.7,0.7);
  \node[font=\bfseries\small, color=red!85] at (2.6,0.1) {$?$};
  \node[font=\bfseries\tiny, color=black!70] at (2.6,-0.35) {Ngày đầu tuần};
\end{tikzpicture}
\vspace{-2.5mm}\end{center}
Sau ngày Chủ nhật thì ngày tiếp theo sẽ lặp lại là ngày nào?
\begin{multicols}{4}
\choice
{Thứ Bảy}
{Thứ Ba}
{\True Thứ Hai}
{Chủ nhật}
\end{multicols}
\loigiai{
Hết ngày Chủ nhật, tuần lễ mới lại bắt đầu và lặp lại từ ngày Thứ Hai.\\
\textbf{Đáp án: C.}
}
\end{ex}

\newpage
% --- TRANG 4: PHẦN 2 - CÂU HỎI ĐÚNG / SAI (4 CÂU) ---

\textbf{Phần 2. Câu trắc nghiệm đúng sai}
\setcounterpageref{ex}{0}

% CÂU 13: Kiến thức tổng quát tuần lễ
\begin{ex} %[1H9N35]
Bảng tổng hợp các ngày trong một tuần lễ:
\begin{center}\vspace{-2.5mm}
\begin{tikzpicture}
  \VeDaiBayNgay{0,0}{0}
\end{tikzpicture}
\vspace{-2.5mm}\end{center}
Xét tính đúng hoặc sai của các khẳng định sau:
\choiceTF
{\True Một tuần lễ có đúng $7$ ngày liên tiếp.}
{Học sinh tiểu học đi học từ Thứ Hai đến Thứ Bảy là $6$ ngày.}
{\True Hai ngày nghỉ cuối tuần là Thứ Bảy và Chủ nhật.}
{\True Ngày đầu tiên trong tuần đi học là Thứ Hai.}
\loigiai{
\begin{itemize}
    \item a) Đúng: Một tuần lễ gồm 7 ngày.
    \item b) Sai: Học sinh đi học 5 ngày (từ Thứ Hai đến Thứ Sáu).
    \item c) Đúng: Cuối tuần nghỉ học Thứ Bảy và Chủ nhật.
    \item d) Đúng: Thứ Hai là ngày đầu tuần.
\end{itemize}
}
\end{ex}

% CÂU 14: Chuỗi ngày liên tiếp (Hôm nay: Thứ Tư)
\begin{ex} %[1H9H35]
Quan sát bộ ba tờ lịch liên tiếp:
\begin{center}\vspace{-2.5mm}
\begin{tikzpicture}
  \VeToLich{-3.4,0}{THỨ BA}{Hôm qua}{myteal}
  \VeToLich{0,0}{THỨ TƯ}{Hôm nay}{myblue}
  \VeToLich{3.4,0}{THỨ NĂM}{Ngày mai}{myorange}
\end{tikzpicture}
\vspace{-2.5mm}\end{center}
Xét tính đúng hoặc sai của các khẳng định sau:
\choiceTF
{\True Nếu hôm nay là Thứ Tư thì ngày mai là Thứ Năm.}
{\True Hôm qua của ngày Thứ Tư là ngày Thứ Ba.}
{Ngày kia (sau hôm nay $2$ ngày) là ngày Thứ Bảy.}
{\True Từ Thứ Tư đến ngày Chủ nhật cùng tuần còn cách $4$ ngày nữa.}
\loigiai{
\begin{itemize}
    \item a) Đúng: Liền sau Thứ Tư là Thứ Năm.
    \item b) Đúng: Liền trước Thứ Tư là Thứ Ba.
    \item c) Sai: Ngày kia là Thứ Sáu (không phải Thứ Bảy).
    \item d) Đúng: Thứ Tư đến Chủ nhật là 4 ngày ($7 - 3 = 4$).
\end{itemize}
}
\end{ex}

% CÂU 15: Thời khóa biểu của Nam
\begin{ex} %[1H9V35]
Thời khóa biểu môn học các ngày trong tuần của Nam:
\begin{center}\vspace{-2.5mm}
\begin{tikzpicture}[scale=0.62]
  \foreach \i/\thu/\mon/\mau in {
    1/{Thứ 2}/{Toán}/myblue,
    2/{Thứ 3}/{Tiếng Việt}/myteal,
    3/{Thứ 4}/{Mĩ thuật}/myorange,
    4/{Thứ 5}/{Âm nhạc}/mypurple,
    5/{Thứ 6}/{Thể dục}/myred
  } {
    \pgfmathsetmacro{\x}{(\i-3)*2.5}
    \draw[fill=\mau!10, draw=\mau!80, line width=0.8pt, rounded corners=3pt] (\x-1.0,-0.6) rectangle (\x+1.0,0.6);
    \node[font=\bfseries\footnotesize, color=\mau!90!black] at (\x,0.18) {\thu};
    \node[font=\bfseries\tiny, color=black!85] at (\x,-0.25) {\mon};
  }
\end{tikzpicture}
\vspace{-2.5mm}\end{center}
Xét tính đúng hoặc sai của các khẳng định sau:
\choiceTF
{\True Nam học môn Mĩ thuật vào ngày Thứ Tư.}
{\True Môn Thể dục diễn ra vào ngày cuối cùng của tuần đi học (Thứ Sáu).}
{Nam học môn Âm nhạc trước môn Mĩ thuật $1$ ngày.}
{\True Nam được nghỉ học ở nhà vào Thứ Bảy và Chủ nhật.}
\loigiai{
\begin{itemize}
    \item a) Đúng: Thứ Tư học Mĩ thuật.
    \item b) Đúng: Thứ Sáu là ngày đi học cuối cùng trong tuần.
    \item c) Sai: Âm nhạc học Thứ Năm (sau Mĩ thuật 1 ngày).
    \item d) Đúng: Cuối tuần Nam được nghỉ học.
\end{itemize}
}
\end{ex}

% CÂU 16: Kế hoạch du lịch của gia đình Rô-bốt
\begin{ex} %[1H9V35]
Lịch trình chuyến du lịch cuối tuần của gia đình Rô-bốt:
\begin{center}\vspace{-2.5mm}
\begin{tikzpicture}
  \VeToLich{-3.4,0}{THỨ SÁU}{Khởi hành}{myblue}
  \VeToLich{0,0}{THỨ BẢY}{Tắm biển}{myorange}
  \VeToLich{3.4,0}{CHỦ NHẬT}{Trở về}{myred}
\end{tikzpicture}
\vspace{-2.5mm}\end{center}
Xét tính đúng hoặc sai của các khẳng định sau:
\choiceTF
{\True Chuyến du lịch diễn ra trong $3$ ngày: Thứ Sáu, Thứ Bảy và Chủ nhật.}
{Gia đình Rô-bốt bắt đầu chuyến đi vào ngày Thứ Năm.}
{\True Ngày kết thúc chuyến du lịch là ngày Chủ nhật.}
{\True Sau khi kết thúc chuyến đi, ngày tiếp theo Rô-bốt đi học trở lại là Thứ Hai.}
\loigiai{
\begin{itemize}
    \item a) Đúng: Chuyến đi gồm 3 ngày: Thứ Sáu, Thứ Bảy, Chủ nhật.
    \item b) Sai: Bắt đầu từ Thứ Sáu (không phải Thứ Năm).
    \item c) Đúng: Trở về vào Chủ nhật.
    \item d) Đúng: Thứ Hai tuần sau Rô-bốt đi học lại.
\end{itemize}
}
\end{ex}

\newpage
% --- TRANG 5: PHẦN 3 - CÂU HỎI TRẢ LỜI NGẮN (4 CÂU) ---

\textbf{Phần 3. Câu hỏi trả lời ngắn}
\setcounterpageref{ex}{0}

% CÂU 17: Một tuần lễ có bao nhiêu ngày?
\begin{ex} %[1H9N35]
Quan sát đầy đủ dải các ngày trong một tuần lễ:
\begin{center}\vspace{-2.5mm}
\begin{tikzpicture}
  \VeDaiBayNgay{0,0}{0}
\end{tikzpicture}
\vspace{-2.5mm}\end{center}
Một tuần lễ có tất cả bao nhiêu ngày?
\textit{Trả lời:} \raisebox{-2pt}{\tcbox[colframe=blue!80!black,colback=white,size=small,arc=3pt]{$7$}}
\loigiai{
Một tuần lễ có $7$ ngày.\\
\textbf{Đáp số: $7$.}
}
\end{ex}

% CÂU 18: Học sinh đi học bao nhiêu ngày trong 1 tuần?
\begin{ex} %[1H9H35]
Em đi học từ ngày Thứ Hai đến hết ngày Thứ Sáu:
\begin{center}\vspace{-2.5mm}
\begin{tikzpicture}[scale=0.7]
  \foreach \i/\thu in {1/{Thứ 2}, 2/{Thứ 3}, 3/{Thứ 4}, 4/{Thứ 5}, 5/{Thứ 6}} {
    \pgfmathsetmacro{\x}{(\i-3)*1.85}
    \draw[fill=blue!12, draw=blue!75, line width=0.8pt, rounded corners=3pt] (\x-0.75,-0.5) rectangle (\x+0.75,0.5);
    \node[font=\bfseries\scriptsize, color=blue!90!black] at (\x,0) {\thu};
  }
\end{tikzpicture}
\vspace{-2.5mm}\end{center}
Trong một tuần lễ, học sinh đi học tất cả bao nhiêu ngày?
\textit{Trả lời:} \raisebox{-2pt}{\tcbox[colframe=blue!80!black,colback=white,size=small,arc=3pt]{$5$}}
\loigiai{
Từ Thứ Hai đến Thứ Sáu gồm có $5$ ngày.\\
\textbf{Đáp số: $5$.}
}
\end{ex}

% CÂU 19: Khoảng cách giữa 2 ngày trong tuần
\begin{ex} %[1H9V35]
Quan sát khoảng thời gian từ Thứ Ba đến Thứ Sáu trong cùng một tuần:
\begin{center}\vspace{-2.5mm}
\begin{tikzpicture}
  \VeToLich{-2.8,0}{THỨ BA}{Bắt đầu}{myteal}
  \VeMuiTenLich{-1.4,0}{2.8}{$? \text{ ngày}$}
  \VeToLich{2.8,0}{THỨ SÁU}{Kết thúc}{myblue}
\end{tikzpicture}
\vspace{-2.5mm}\end{center}
Từ Thứ Ba đến Thứ Sáu trong cùng một tuần cách nhau mấy ngày?
\textit{Trả lời:} \raisebox{-2pt}{\tcbox[colframe=blue!80!black,colback=white,size=small,arc=3pt]{$3$}}
\loigiai{
Khoảng cách thời gian là: $6 - 3 = 3$ ngày (Thứ Tư, Thứ Năm, Thứ Sáu).\\
\textbf{Đáp số: $3$.}
}
\end{ex}

% CÂU 20: Hai tuần lễ có bao nhiêu ngày?
\begin{ex} %[1H9C35]
Bạn An được nghỉ hè trọn vẹn $2$ tuần lễ:
\begin{center}\vspace{-2.5mm}
\begin{tikzpicture}
  \draw[fill=blue!10, draw=blue!75, line width=0.8pt, rounded corners=3pt] (-3.8,-0.55) rectangle (-0.6,0.55);
  \node[font=\bfseries\footnotesize, color=blue!90!black] at (-2.2,0.12) {Tuần $1$};
  \node[font=\bfseries\tiny, color=black!70] at (-2.2,-0.25) {$7$ ngày};
  
  \node[font=\bfseries\Large, color=red!85] at (0,0) {$+$};
  
  \draw[fill=orange!10, draw=orange!75, line width=0.8pt, rounded corners=3pt] (0.6,-0.55) rectangle (3.8,0.55);
  \node[font=\bfseries\footnotesize, color=orange!90!black] at (2.2,0.12) {Tuần $2$};
  \node[font=\bfseries\tiny, color=black!70] at (2.2,-0.25) {$7$ ngày};
\end{tikzpicture}
\vspace{-2.5mm}\end{center}
Hỏi bạn An được nghỉ hè tất cả bao nhiêu ngày?
\textit{Trả lời:} \raisebox{-2pt}{\tcbox[colframe=blue!80!black,colback=white,size=small,arc=3pt]{$14$}}
\loigiai{
Mỗi tuần lễ có $7$ ngày. Hai tuần lễ có: $7 + 7 = 14$ ngày.\\
\textbf{Đáp số: $14$.}
}
\end{ex}
